\section{Discussion}\label{sec:discussion}

\paragraph{When is the reduction tight?} Optimizing \cref{thm:reduction} over descriptions gives, up to estimation error,
\[
\err(f)\ \lesssim\ \min_\sigma\Big[\err(g_\sigma)+\Pr_{D_X}[f\ne g_\sigma]+\tfrac{L(\sigma)\ln 2}{m}\Big]\ \le\ \err(f)+2\epsilon+\tfrac{\ln 2}{m}\,K_\epsilon(f),
\]
where $K_\epsilon(f)=\min\{L(\sigma):\Pr_{D_X}[f\ne g_\sigma]\le\epsilon\}$ is the $\epsilon$-approximate description length of $f$ under $D_X$ (the first inequality uses $\err(g)\le\err(f)+\epsilon$; the rough form uses the realizable rate $L\ln2/m$, whereas the binomial inversion adds a term of order $\sqrt{\err(g)L/m}$ when $g$ has training errors, so it is optimistic there). So the reduction is as good as $f$'s \emph{lossy} compressibility under the input distribution. Lossy is essential. A network that interpolates noisy labels differs from the clean target only near the training points, a set of small $D_X$-measure, so it can have small $K_\epsilon$. In our label-noise experiments the bound on such interpolating networks is 0.17 against a test error of 0.14 (\cref{sec:exp-bool}). The bound fails, correctly, when $f$ is far from every short program; this includes the case of a good network for an incompressible target. Generalization of a learning \emph{algorithm} can always be reduced to compression in the other sense, that of the sample (Littlestone--Warmuth; \citet{moran2016sample}), with the training procedure as decoder. This is a second description language for the same reduction, and its decoder can use the initialization for free.

\paragraph{Provable closeness without inputs from $D_X$.} The second term of \cref{thm:reduction} is where a validation-free bound has to work. With known or sampleable $D_X$ it is free; with unlabeled inputs it is the method of \citet{bazinet2026bound}. With neither, the only information about $D_X$ is the training inputs, which $f$ depends on. The dithered transfer theorem resolves this by making the certificate a function of the description alone. It costs a weight-space description of $f$ precise enough for $f$'s margins, and the cost grows with the number of weights that matter at typical inputs. Functional simplicity of $f$ (e.g.\ $f$ is close to a 5-junta) is invisible to it, unlike to \cref{thm:reduction}. Closing this gap is the main open problem the experiments point to. A certificate would have to establish $f(x)=g(x)$ at a point from a description of $g$ plus a short proof about $f$, as in the compact-proofs programme of \citet{gross2024compact}, rather than from weight-space proximity.

\paragraph{The initialization.} The worry that the random initialization must be described is unfounded at the level of validity (\cref{lem:init}). Any data-independent randomness can be used by the code for free, and a seed costs 32 bits anyway. What remains is a question about the update. (i) Function-space codecs (our Boolean truth tables and distilled students) never refer to $f$'s initialization, so it cannot cost anything. (ii) Weight-space codecs, which training-data-only bounds need, code the update relative to $\theta_0$. Weights that receive no gradient (first-layer weights on pixels that are zero in every training image) have update exactly zero and cost almost nothing, whereas coding them relative to zero costs their full entropy. (iii) Exact zero initialization is a degenerate point for ReLU MLPs (all hidden units are dead and only the output bias receives a gradient), so it does not train. A deterministic structured initialization trains and costs nothing, but a seeded random one also costs nothing. (iv) From a typical initialization, wide networks in the mean-field scaling (or with a centred initialization) trained by gradient descent approach a deterministic limit, so the initialization only adds fluctuations of order $\mathrm{width}^{-1/2}$. In the NTK scaling the trained function is not deterministic: it is the kernel regression plus the initial function, an $O(1)$ Gaussian-process sample. A function-space codec absorbs such fluctuations into the disagreement term; a weight-space codec must pay for them or certify through them. In the linearized regime the free eigenbasis of the initial kernel gives a codec (\cref{prop:kernel}) whose cost, heuristically, grows like the $n^k$ sample complexity of kernel methods on degree-$k$ targets; we observed the predicted rank $1+n$ for majority targets (signs of degree-1 functions) at one value of $n$. The rich regime, in which the network finds the relevant coordinates, is what makes trained networks short to describe.

\paragraph{Limitations.} The training-data-only bound is proved for one-hidden-layer networks. The same argument extends layer by layer (pre-activation errors stay sub-Gaussian, though no longer independent across units), but we have not carried it out. Our codecs and step grids were designed with the datasets in mind, as in all of this literature: strictly, the code must be fixed before the sample is drawn. The Boolean experiments use small input dimension and synthetic targets. On images we bound MLPs, not convolutional networks, and our unlabeled-data bounds are far from the best published ones \cite{bazinet2026bound}, which use data-dependent priors and stronger compressors.
